\documentclass[10pt,twocolumn]{article}
\usepackage[T1]{fontenc}
\usepackage[utf8]{inputenc}
\usepackage{lmodern}
\usepackage[margin=1.8cm,top=2.4cm,bottom=2cm,columnsep=0.6cm]{geometry}
\usepackage{fancyhdr}
\usepackage{titlesec}
\usepackage{booktabs}
\usepackage{amsmath,amssymb,amsfonts}
\usepackage{microtype}
\usepackage{xcolor}
\usepackage{mdframed}
\usepackage{parskip}
\usepackage{enumitem}
\usepackage{hyperref}

\definecolor{rulered}{RGB}{180,30,30}
\definecolor{boxgray}{RGB}{245,245,245}
\definecolor{keywordgray}{RGB}{100,100,100}

\pagestyle{fancy}
\fancyhf{}
\fancyhead[L]{\textsc{\small Claude Mag}}
\fancyhead[C]{\small\textit{Machine Learning Research Digest}}
\fancyhead[R]{\small 2026-08-19}
\fancyfoot[C]{\small\thepage}
\renewcommand{\headrulewidth}{0.8pt}

\titleformat{\section}{\normalsize\bfseries\color{rulered}}{}{0pt}{}%
  [\vspace{-4pt}\color{rulered}\rule{\columnwidth}{0.4pt}]
\titlespacing{\section}{0pt}{8pt}{4pt}

\hypersetup{colorlinks=true,urlcolor=rulered,linkcolor=black}

\begin{document}
\setlength{\parindent}{0pt}
\setlength{\parskip}{4pt}

{\small\color{keywordgray}\textbf{Keywords:}
  representational geometry \textbullet{}
  classification \textbullet{}
  prototype theory \textbullet{}
  universal representations}\par\vspace{4pt}

{\LARGE\bfseries\raggedright
  Sparse Prototype Code Underlies Classification
  and Prediction Across Modalities}\par\vspace{4pt}

{\large\itshape\raggedright
  A Universal Geometry of Within-Class Variability, Aligned With Class
  Centroids and Competing Classes, Emerges Across Vision, Audio,
  and Language Models}\par\vspace{6pt}

{\small\textbf{By} Yehonatan Avidan, Daniel D.\ Lee, Haim Sompolinsky}\par\vspace{2pt}

{\small\color{keywordgray}arXiv:2608.15632 \textbullet{} Preprint, August 2026}
\par\vspace{2pt}\noindent{\color{rulered}\rule{\columnwidth}{0.6pt}}\vspace{6pt}

Why do deep networks trained on classification tasks learn representations that
generalize so broadly? This paper identifies a universal representational geometry
--- the Sparse Prototype Code --- that appears across state-of-the-art models in
vision, audio, and language. The key discovery: within-class variability in
representation space is not random. Its classifier-relevant component aligns
strongly with the class prototype centroid and with the centroids of a sparse
set of competing classes. This structure explains inter-class confusion patterns,
predicts transferability, and provides a unified framework for understanding why
classification training produces broadly useful representations.

\begin{mdframed}[backgroundcolor=boxgray,linecolor=rulered,linewidth=1pt,
    innerleftmargin=6pt,innerrightmargin=6pt,
    innertopmargin=6pt,innerbottommargin=6pt]
\textbf{\small AT A GLANCE}\par\vspace{4pt}
\begin{description}[leftmargin=0pt,labelsep=4pt,itemsep=2pt,parsep=0pt,
    font=\small\bfseries]
  \item[Problem:] The geometric structure of learned representations is poorly
    understood, limiting principled model design and analysis.
  \item[Why it matters:] A universal geometry across modalities would explain
    why classification training generalizes and enable principled comparison
    of representations.
  \item[Method:] Spectral analysis of within-class covariance matrices in
    final representation spaces of vision, audio, and language models.
  \item[Result:] Sparse Prototype Code geometry confirmed across all tested
    architectures and modalities; predicts confusion patterns and transferability.
  \item[Evidence:] 33 pages, 13 figures, 14 tables spanning ViT, ResNet,
    CLIP, wav2vec, BERT, and GPT-class models.
\end{description}
\end{mdframed}

\section{The Big Picture}

Deep learning models trained on classification tasks have proven surprisingly
versatile: representations learned to classify ImageNet transfer to medical imaging,
satellite imagery, and fine-grained recognition; representations learned on
language classification transfer to generation, translation, and reasoning.
The mechanism behind this broad transferability is not well understood.

The dominant hypothesis is that classification forces networks to learn
disentangled, semantically meaningful features. But this is vague: what
geometric structure do these features have in representation space? This paper
provides a precise characterization: the Sparse Prototype Code. It shows
that this structure is not specific to one modality or architecture but is
a universal consequence of training on classification tasks.

\section{Why Existing Approaches Fall Short}

\textbf{Neural collapse theory} (Papyan et al., 2020) characterizes the
geometry of the final layer at the end of training: class means collapse
to the vertices of a simplex equiangular tight frame (ETF). This is a
useful characterization of inter-class geometry but does not address
\emph{within-class} variability — the spread of examples around each class
mean. Neural collapse predicts that within-class variability vanishes at
convergence, which is not the case for practical models trained on realistic
data.

\textbf{CKA and representation similarity} methods compare representations
between models but do not characterize the intrinsic geometry of a single
model's representation space.

\textbf{Prototype theory} in cognitive science posits that categories are
organized around prototypical examples, but does not make quantitative
predictions about the geometric structure of computational model representations.

The Sparse Prototype Code bridges these gaps by providing a quantitative,
testable geometric characterization of within-class variability.

\section{Core Mechanism}

For a model $f_\theta$ and class $c$, let $\mu_c = \mathbb{E}_{x \sim p(x|c)}[f_\theta(x)]$
be the class centroid and $\Sigma_c = \mathbb{E}_{x \sim p(x|c)}[(f_\theta(x) - \mu_c)(f_\theta(x) - \mu_c)^\top]$
be the within-class covariance.

The paper performs spectral analysis of $\Sigma_c$, decomposing it as
$\Sigma_c = U_c \Lambda_c U_c^\top$. The key finding is that the leading
eigenvectors $u_1, u_2, \ldots$ of $\Sigma_c$ align strongly with class
centroids:
\begin{enumerate}[itemsep=2pt,parsep=0pt]
  \item $u_1 \approx \mu_c / \|\mu_c\|$ (the class prototype direction)
  \item $u_2, \ldots, u_k \approx \mu_{c_j} / \|\mu_{c_j}\|$ for a sparse
    set of competing classes $\{c_j\}_{j=1}^{k}$, with $k \ll C$ (where
    $C$ is the total number of classes)
\end{enumerate}

The eigenvalue spectrum is sparse: a small number of eigenvalues are large
(aligned with prototype and competing-class directions) and the remaining
eigenvalues are small (random variability). This sparsity is the defining
feature of the Sparse Prototype Code.

\section{Technical Formulation}

Formally, define the \emph{prototype alignment} of class $c$:
\[
  A_{\text{proto}}(c) = \frac{u_1^\top \mu_c}{\|\mu_c\|}
\]
and the \emph{competing-class alignment} of class $c$ with class $c'$:
\[
  A_{\text{comp}}(c, c') = \max_{j \geq 2} \frac{u_j^\top \mu_{c'}}{\|\mu_{c'}\|}
\]

The Sparse Prototype Code hypothesis states:
\[
  A_{\text{proto}}(c) \approx 1 \quad \text{for all } c
\]
\[
  \sum_{c' \neq c} \mathbf{1}[A_{\text{comp}}(c,c') > \tau] = k(c) \ll C
\]
where $k(c)$ is a sparse count of strongly competing classes.

The paper verifies both conditions across all tested models and modalities.

\section{Learning or Inference Procedure}

The analysis is purely empirical and requires no training. For each model:
\begin{enumerate}[itemsep=2pt,parsep=0pt]
  \item Extract final-layer representations for all examples in the test set.
  \item Compute class centroids $\{\mu_c\}$.
  \item Compute within-class covariance matrices $\{\Sigma_c\}$.
  \item Perform eigendecomposition of each $\Sigma_c$.
  \item Measure alignment of leading eigenvectors with class centroids.
\end{enumerate}
All steps are straightforward linear algebra; no optimization is involved.

\section{What the Guarantee Says}

The paper provides a theoretical derivation showing that the Sparse Prototype
Code is the \emph{optimal} within-class geometry for a linear classifier that
minimizes cross-entropy loss while maintaining within-class coherence. That is,
gradient descent on classification loss with weight decay provably induces the
Sparse Prototype Code in the final layer in the limit of large data and a
sufficiently expressive model. This is not a finite-sample guarantee but
provides a theoretical grounding for the empirical observations.

\section{Experimental Findings}

\begin{itemize}[itemsep=2pt,parsep=0pt]
  \item \textbf{Vision (ViT, ResNet, CLIP):} Prototype alignment $A_{\text{proto}} > 0.95$
    for all classes in all models. Average competing-class count $k(c) = 3$--$8$
    out of 1000 classes (ImageNet). Competing classes match human-annotated
    confusable categories.
  \item \textbf{Audio (wav2vec):} Same geometry emerges in the audio representation
    space for AudioSet classification. Competing classes are acoustically similar
    sounds (e.g., ``bark'' and ``growl'').
  \item \textbf{Language (BERT, GPT-class):} Sparse Prototype Code confirmed in
    text classification tasks. Competing classes correspond to semantically
    adjacent categories (e.g., ``positive'' and ``mildly positive'' sentiment).
  \item \textbf{Transferability prediction:} Models with higher prototype alignment
    and lower competing-class count transfer better to new tasks (Pearson $r = 0.78$).
\end{itemize}

\section{Ablations and Interpretation}

\textbf{Does the geometry require classification training?} Models trained with
contrastive losses (CLIP, SimCLR) show partial but weaker Sparse Prototype Code
structure, with higher competing-class count and lower prototype alignment.
Supervised classification induces stronger structure.

\textbf{Does architecture matter?} The geometry is consistent across ViT and
ResNet at matched model capacity, suggesting it is a property of the training
objective rather than the architecture.

\textbf{Depth:} The Sparse Prototype Code is weakest in early layers and
strongest in the final representation layer, consistent with hierarchical
abstraction during training.

\section*{Reference}

Yehonatan Avidan, Daniel D.\ Lee, Haim Sompolinsky.
\textbf{Sparse Prototype Code Underlies Classification and Prediction Across Modalities}.
arXiv:2608.15632, August 2026.
\url{https://arxiv.org/abs/2608.15632}

\end{document}
